\section{Ratios de solvencia}
La solvencia analiza la estructura de financiamiento. Una empresa puede tener liquidez de corto plazo y, aun así, presentar riesgo si su deuda total es elevada o si los gastos financieros presionan la utilidad.

\subsection{Razón de endeudamiento}
\[
\text{Endeudamiento}=\frac{\text{Pasivo total}}{\text{Activo total}}\times 100
\]
\textbf{Caso académico elaborado con fines educativos.} Industrial Callao S.A.C. registra activos por \soles{2,000,000.00} y pasivos por \soles{900,000.00}.
\[
\frac{900,000.00}{2,000,000.00}\times 100=45.00\text{ \char37}
\]
\textbf{Interpretación:} el 45.00 \char37{} de los activos ha sido financiado mediante obligaciones con terceros.
\textbf{Decisión:} evaluar costo financiero, vencimientos y capacidad de pago antes de aumentar deuda.

\subsection{Relación deuda-patrimonio}
\[
\text{Deuda-patrimonio}=\frac{\text{Pasivo total}}{\text{Patrimonio}}
\]
\[
\frac{900,000.00}{1,100,000.00}=0.82
\]
\textbf{Interpretación:} por cada \soles{1.00} aportado o acumulado en patrimonio, existen \soles{0.82} financiados por terceros.
